% ЕМТ 2024, X клас — точен LaTeX препис от официалните файлове в Klasirane.com.
% Условия: https://klasirane.com/api/competitions/EMT/All/2024/10%20%D0%BA%D0%BB/probs
% SHA-256 (условия): 48d43c0f91a46b1c524260587eec3c0ed997f7ce7feebe3c06dbdaf40405d061
% Решения: https://klasirane.com/api/competitions/EMT/All/2024/10%20%D0%BA%D0%BB/sol
% SHA-256 (решения): 77561cf8e68f23f94747d79fdde8491991a7e9956e34f0e4bbc89de7def79fd1
% Точковите оценки, правописът, формулировките и видимите печатни особености са запазени от източника.
% И двата PDF файла съдържат сканирани страници; видимите страници са проверени пряко.
% В решение 10.1 двете квадратни уравнения във втория метод са преписани с отпечатаната липса на множителя t.
% В решение 10.2 равенството с BL'P е преписано точно както е отпечатано, включително липсващия знак за ъгъл.

Задача 10.1. Да се намерят всички двойки реални числа $(x;y)$, които са решения на системата
$$
\begin{cases}
(x^2+xy+y^2)\sqrt{x^2+y^2}=88,\\
(x^2-xy+y^2)\sqrt{x^2+y^2}=40.
\end{cases}
$$

Отговор. 4 решения: $(x;y)=(\sqrt7\pm1,\sqrt7\mp1); (x;y)=(-\sqrt7\pm1,-\sqrt7\mp1)$.

Решение. Първи метод. Тъй като $x^2+y^2\geq0$, то в дефиниционната област на квадратния корен попадат всички реални $(x;y)$. Събирайки двете уравнения и разделяйки всяка от двете страни на две, получаваме
$$
\left(\sqrt{x^2+y^2}\right)^3=64=4^3\quad\Longrightarrow\quad \sqrt{x^2+y^2}=4\text{ и }x^2+y^2=16.
$$
Изваждайки от първото уравнение второто и разделяйки всяка от двете страни на две, получаваме
$$
xy\sqrt{x^2+y^2}=24\quad\Longrightarrow\quad xy=24/4=6\quad\Longrightarrow\quad x^2y^2=36.
$$
Според формулите на Виет $x^2$ и $y^2$ са корените на уравнението
$$
t^2-16t+36=0\quad\Longrightarrow\quad t_{1,2}=8\pm2\sqrt7=(\sqrt7\pm1)^2.
$$
Оттук $x,y\in\{\sqrt7\pm1,-\sqrt7\pm1\}$. От съображения за симетрия виждаме, че ако $(x;y)$ е решение, то решения са и $(y;x)$, $(-x;-y)$ и $(-y;-x)$. Тъй като $xy=6>0$, то $x$ и $y$ са с еднакви знаци. Следователно всяка една от четирите възможности за $x\in\{\sqrt7\pm1,-\sqrt7\pm1\}$ еднозначно определя съответното $y$ и така получаваме четирите решения на задачата $(x;y)=(\sqrt7\pm1,\sqrt7\mp1); (x;y)=(-\sqrt7\pm1,-\sqrt7\mp1)$.

Втори метод. От $x^2+y^2=16$ и $xy=6$ получаваме $(x+y)^2=16+2\cdot6=28$, откъдето $|x+y|=2\sqrt7$. Ако $x+y=2\sqrt7$, то $x$ и $y$ са корени на квадратното уравнение
$$
t^2-2\sqrt7+6=0\quad\Longleftrightarrow\quad t_{1,2}=\sqrt7\pm1,
$$
и поради симетрията на $x$ и $y$ в този случай уравнението има 2 решения: $(x;y)=(\sqrt7\pm1;\sqrt7\mp1)$.

Ако пък $x+y=-2\sqrt7$, то $x$ и $y$ са корени на квадратното уравнение
$$
t^2+2\sqrt7+6=0\quad\Longleftrightarrow\quad t_{1,2}=-\sqrt7\pm1,
$$
и поради симетрията на $x$ и $y$ в този случай уравнението има 2 решения: $(x;y)=(-\sqrt7\pm1;-\sqrt7\mp1)$.

Оценяване. (6 точки) Първи метод: По 1 т. за $x^2+y^2=16$ и $xy=6$; 2 т. за $x,y\in\{\sqrt7\pm1,-\sqrt7\pm1\}$; 2 т. за намирането на четирите решения. Втори метод: по 1 т. за $|x+y|=2\sqrt7$ и $xy=6$; по 2 т. за всеки от случаите $x+y=\pm2\sqrt7$, минус 1 т., ако е изпусната симетрията $(x;y)\to(y;x)$. И при двата подхода, ако преобразуванията не са еквивалентни, а само следствени (т.е., $\Longleftrightarrow$ е заменено с $\Longrightarrow$), то се отнема 1 т. при липса на проверка.

Задача 10.2. Даден е неравнобедрен остроъгълен триъгълник $ABC$, в който $AL$ ($L\in BC$) е ъглополовящата на $\angle BAC$ и $M$ е средата на $BC$. Нека ъглополовящите на $\angle AMB$ и $\angle CMA$ пресичат $AB$ и $AC$ съответно в точките $P$ и $Q$. Да се докаже, че описаната окръжност около триъгълника $APQ$ се допира до $BC$ тогава и само тогава, когато минава през $L$.

Решение. От свойството на ъглополовящата и факта, че $BM=MC$, получаваме
$$
\frac{AP}{PB}=\frac{AM}{BM}=\frac{AM}{MC}=\frac{AQ}{QC}
$$
и от теоремата на Талес $PQ\parallel BC$. Следователно $\angle BLP=\angle LPQ=\varphi$. Нека $\omega$ е описаната окръжност около триъгълника $APQ$.

Нека $\omega$ се допира до $BC$ в точка $L'$. Тогава
$$
\angle L'QP=BL'P=\angle L'PQ,
$$
т.е. $L'$ е средата на дъгата $\overset{\frown}{PQ}$ от $\omega$. Следователно $AL'$ е ъглополовящата на $\angle PAQ$, т.е. $L'\equiv L$.

Нека $L\in\omega$. Тогава понеже $AL$ е ъглополовяща на $\angle PAQ$, то $PL=LQ$, откъдето
$$
\angle PQL=\angle LPQ=\angle BLP=\varphi.
$$
Последното означава, че $\omega$ се допира до $BC$.

Оценяване. (6 точки) 2 т. за $PQ\parallel BC$, по 2 т. за всяка от двете посоки.

Задача 10.3. Да се намерят всички полиноми $P$ с цели коефициенти, за които съществува естествено число $N$, такова че за всяко $n\geq N$, всеки прост делител на $n+2^{\lfloor\sqrt n\rfloor}$ е делител и на $P(n)$. (Тук $\lfloor x\rfloor$ означава най-голямото цяло число, по-малко или равно на реалното число $x$.)

Отговор. $P\equiv0$.

Решение. Очевидно $P\equiv0$ върши работа. Ще покажем, че други $P$ няма. Да забележим, че $\lfloor\sqrt n\rfloor=k$ за всеки две естествени $n$ и $k$, за които $n\in[k^2;k^2+2k]$. Да фиксираме естествено число $m\geq N$ и простото число $p$. Да положим в горното $k=p-1+m$ и да разгледаме интервала $[k^2;k^2+2k]$. Той има дължина $2k>p$, следователно в него съществува естествено число $n\equiv-2^m\pmod p$. Тогава от теоремата на Ферма получаваме
$$
n+2^{\lfloor\sqrt n\rfloor}=n+2^k\equiv-2^m+2^{p-1+m}\equiv0\pmod p.
$$
Следователно $p\mid P(n)$. Понеже $P$ е полином с цели коефициенти,
$$
P(-2^m)\equiv P(n)\equiv0\pmod p.
$$
Така $p\mid P(-2^m)$ за всяко просто число $p$. Следователно $P(-2^m)=0$ за всяко естествено число $m\geq N$. Това значи, че $P\equiv0$, с което решението е завършено.

Оценяване. (7 точки) 1 т. за ясната идея за разглеждане на числа $n$ в интервали от вида $[k^2;k^2+2k]$, 4 т. за конструиране на двойки $(n;p)$ от посочения вид и доказване на делимостта, 2 т. за довършване.

Задача 10.4. Дадени са пълен ориентиран граф $G$ с 2024 върха и естествено число $k\leq10^5$. Ангел и Борис играят следната игра: Ангел оцветява $k$ от ребрата на $G$ в червено и поставя пул в един от върховете на $G$. След това двамата правят ходове, редувайки се; Ангел започва. На всеки свой ход той премества пула в съседен връх, след което Борис променя ориентацията на някое от ребрата, което не е червено. Ако в някакъв момент Ангел не може да премести пула, то той губи и победител е Борис. Да се определи в зависимост от $G$ и $k$ дали Борис има печеливша стратегия.

Отговор. Борис има печеливша стратегия при $k\leq2$ и всякакъв граф $G$, както и при $k\geq3$, когато в $G$ няма цикли.

Решение. Ще започнем решението със следните две леми.

Лема 1. Нека $G$ е ориентиран граф с $n$ върха, в който няма цикли. Тогава можем да номерираме върховете му с естествените числа от 1 до $n$ така, че ако за някои два върха $x$ и $y$ с номера $i$ и $j$ съответно има ребро $x\to y$, то непременно $i<j$.

Решение. Ще използваме индукция по $n$. При $n=2$ твърдението е тривиално. Нека сега исканото е в сила за всеки граф с по-малко от $n$ върха. Ако допуснем, че от всеки връх в $G$ излиза поне по едно ребро, то бихме получили цикъл, което е противоречие. Значи съществува връх $v$, от който не излизат ребра. Номерираме него с $n$, а останалите върхове номерираме съгласно индуктивната хипотеза за графа $G_1:=G\setminus\{v\}$ (понеже в $G$ няма цикли, то и в $G_1$ няма). Това гарантира, че условието за ребрата в $G_1$ се изпълнява. Останалите ребра са от вида $u\to v$, но номерът на $v$ е най-големият възможен ($n$), значи и за тях исканото е в сила. С това лемата е доказана.

Лема 2. Нека $G$ е пълен ориентиран граф с $n$ върха, в който има поне един цикъл. Тогава в $G$ има цикъл с дължина 3.

Решение. Да допуснем противното и да разгледаме цикъл $v_1v_2\ldots v_kv_1$ с минимална дължина $k\geq4$. Ако има ребро $v_3\to v_1$, то $v_1v_2v_3v_1$ е цикъл с дължина $3<k$, което е противоречие. Значи имаме реброто $v_1\to v_3$; сега обаче $v_1v_3\ldots v_kv_1$ е цикъл с дължина $k-1<k$, което отново е противоречие. Следователно в $G$ има цикъл с дължина 3, с което лемата е доказана.

Сега да пристъпим към решението. Ако $k\geq3$ и в $G$ има цикъл, то съгласно Лема 2 в $G$ има цикъл $C$ с дължина 3. Значи Ангел може да оцвети $C$ и още произволни $k-3$ ребра в червено и да постави пула в един от върховете на $C$. Понеже Борис не може да промени ориентацията на ребрата на този цикъл, то Ангел може да направи неограничен брой ходове, което го прави победител.

Нека сега е вярно, че в $G$ или няма цикъл, или $k\leq2$. Ще докажем, че и в двата случая Борис може след краен брой промени (възможно нула) да приведе $G$ в граф, който може да бъде номериран със свойството в Лема 1. В първия случай това следва директно от лемата, остава да го проверим за $k\leq2$. Да разгледаме червените ребра, които Ангел оцветява. Понеже те двете не образуват цикъл, то можем да номерираме участващите в тях върхове по желания начин. Останалите върхове в $G$ номерираме по произволен начин и докато има ребра от вида $i\to j$, където $i>j$, Борис променя тяхната ориентация. Така получихме пълен ориентиран граф $H$ с върхове числата от 1 до 2024 и ребро $i\to j$ за всеки $1\leq i<j\leq2024$. Нека след преместването от страна на Ангел пулът се намира във връх $n$. Разделяме върховете на $H$ на 2 групи по следния начин:
$$
A=\{1,2,\ldots,1000\};\quad B=\{1001,1002,\ldots,2024\}.
$$
Да забележим, че ако ребрата, излизащи от върха, в който се намира пулът (в случая се намира в $n$) преди хода на Ангел, са само към такива с по-голям номер, то непременно върхът, в който пулът бива преместен, има по-голям номер. Остава да покажем, че Борис може да поддържа това свойство в сила.

Ребрата между двойка върхове от една и съща група са повече от $10^5$ и за двете групи, следователно и в двете групи има поне по едно ребро, което не е червено – нека това са $a=a_1\to a_2$ (респективно от върхове в $A$) и $b=b_1\to b_2$ (респективно от върхове в $B$). Действаме по следния начин:

• ако пулът е във връх от $A$, променяме $b$;

• ако пулът е във връх от $B\setminus\{b_2\}$, променяме $a$;

• ако пулът е във $b_2$ и имаме реброто $b_1\to b_2$, променяме $a$;

• ако пулът е във $b_2$ и имаме реброто $b_2\to b_1$, променяме $b$.

Лесно се съобразява, че исканото се изпълнява при тази стратегия.

Оценяване. (7 точки) 3 т. за пълно описание на случая, в който Ангел печели, 3 т. за пълно описание на случая, в който Борис печели, 1 т. за завършване. При формулиране на стратегията на Ангел според дължината на минималния цикъл на $G$ и отстъпване на Лема 2 се отнемат 2 т.
